\paragraph{KT-HIGH-ORDER-TWIST-TARGET-SCOPE.}
The high-order twists used in Heumann--Vatsal Theorem~5.1
cannot simply replace the quadratic auxiliary field in the present
argument. This is a scope statement about the displayed source
specialization, not an impossibility theorem for a different construction.

The exact declaration \texttt{UpperTheory.V.HV.GD} in
\texttt{HeightAndVanishing/Counts.lean}, lines~9--22, requires
the field to have degree two and be totally complex. Its condition
\texttt{ine} says that the \emph{component at the coefficient prime}
of the finite character has square one. It does not say that the
entire finite character is quadratic: the definition of
\texttt{comp} is in
\texttt{HeightAndVanishing/DirichletPrimeToLPartMulLPowerPart.lean},
lines~75--86. High-order characters away from that prime are
therefore not excluded just by this particular condition.

They do, however, change the target. In the latter file,
\texttt{Vlen} at line~116 is
\texttt{Chr.V'} evaluated at the Galois character
\texttt{Chr.dirG} of the actual finite character. The orbit in
lines~140--145 consists of characters with one fixed kernel.
A nontrivial character does not belong to the trivial character's
orbit. The rational untwisted target used here is the
specialization with finite character one.

There is also no direct reduction-to-one substitution using a
nontrivial character of order prime to $p$. Reduction is injective
on roots of unity of order prime to $p$: if $\zeta\ne1$ has
order $m$ prime to $p$ and reduces to one, then
$1+\zeta+\cdots+\zeta^{m-1}=0$ would reduce to $m=0$,
a contradiction. The condition \texttt{GD.ord} expressly
requires prime-to-$p$ order. Thus the high-order character
nonvanishing theorem supplies neither a quadratic field nor a
vanishing result for the original finite-character-one Selmer
group through these declarations alone. An additional comparison
would have to be proved.
